\section{Introduction}

A finite \'etale scheme \(X\) over \(\Q\) can satisfy
\begin{equation}
 X(\R)\ne\varnothing,\qquad X(\Q_p)\ne\varnothing\quad\text{for every prime }p,
 \qquad X(\Q)=\varnothing.
 \label{eq:hasse-condition}
\end{equation}
The classical example of Berend and Bilu
\cite{BerendBilu} is the root scheme of
\begin{equation}
 f_{19}(X)=(X^3-19)(X^2+X+1).
 \label{eq:berend-bilu}
\end{equation}
Its factors cover different completions: the quadratic supplies a root at
primes \(1\pmod3\), while the cubic supplies one at primes \(2\pmod3\)
and at \(3\). Neither factor has a rational root.

The prescribed-fiber construction of \cite{vanRijnFibers} realizes each
finite \'etale algebra of rank at least three as a full fiber of a
Jacobian-one map of affine three-space. Here the map is fixed while the
arithmetic algebra varies. The family
\begin{equation}
 f_a(X)=(X^3-a)(X^2+X+1)
 \label{eq:intersective-family}
\end{equation}
fits one target line because the translation \(X=S-\tfrac12\) centers the
quadratic:
\begin{equation}
 X^2+X+1=S^2+\frac34.
 \label{eq:centered-quadratic-intro}
\end{equation}
Consequently varying \(a\) moves only the quadratic and constant
coefficients. These are the two free coefficients of the quadratic-gauge
inverse equation
\begin{equation}
 E_{\Pi,B,C}(S)=G_\Pi(S)-\frac{g_1}{2}(BS^2+C).
 \label{eq:inverse-intro}
\end{equation}

For every \(a\in\Q\), put
\begin{equation}
 y_a=\left(-1,\frac{32a}{9},\frac{8a+1}{3}\right),
 \qquad \ell=\{y_a:a\in\Q\}.
 \label{eq:target-family}
\end{equation}
Define the rational parameter set
\begin{equation}
 \mathcal H=\left\{a\in\Q^\times\setminus\Q^{\times3}:
 \begin{array}{l}
 v_p(a)\equiv0\pmod3\quad(p\equiv2\pmod3),\\
 a\in\Q_3^{\times3}
 \end{array}\right\}.
 \label{eq:rational-family}
\end{equation}
The \(3\)-adic condition means that \(v_3(a)\) is divisible by three and
the remaining unit is \(1\) or \(-1\pmod9\).

\begin{theorem}[Fixed map and complete rational classification]
\label{thm:main}
There is an explicit polynomial map
\[
 \Phi:\A^3_\Q\longrightarrow\A^3_\Q,
 \qquad \det D\Phi=1,\qquad \gdeg(\Phi)=5,
\]
such that, for every \(a\in\Q\setminus\{0,1\}\),
\begin{equation}
 \Phi^{-1}(y_a)\simeq
 \Spec\Q[X]/\bigl((X^3-a)(X^2+X+1)\bigr).
 \label{eq:main-fiber}
\end{equation}
This is a full reduced fiber of rank five. It is everywhere locally
soluble and has no rational point if and only if \(a\in\mathcal H\).
The exceptional parameters \(0,1\) give no further Hasse failures.
\end{theorem}

The geometric assertion concerns an ordinary polynomial endomorphism of
affine three-space. Since its geometric degree is five, it asserts a
noninvertible Keller map and thus has a negative consequence for the
Jacobian conjecture in dimension three. This is a consequence of the
fixed-map construction, logically separate from the rational local and
counting arguments for the displayed root schemes.

For example, \(19\), \(91\), \(152=2^3\cdot19\), and \(19/37\) belong
to \(\mathcal H\). The latter two illustrate why a sufficient family of
integers supported on primes \(1\pmod3\) does not exhaust the line.

We use the height induced by \((u,v,w)\mapsto[1:u:v:w]\) in
\(\mathbb P^3(\Q)\). For a reduced fraction \(a=m/n\), with \(n>0\),
the exact normalization is
\begin{equation}
 H(y_{m/n})=
 \frac{\max\{9n,32|m|,|24m+3n|\}}
 {\gcd(n,8)\gcd(|m|,3)}.
 \label{eq:height-intro}
\end{equation}
Both the projective content and the last linear form affect the leading
constant. In particular, the height region for positive parameters has a
corner cut from its enclosing rectangle.

\begin{theorem}[All rational parameters on the line]
\label{thm:rational-count-intro}
Let \(N_{\Phi,\ell}(B)\) count the Hasse-failing targets on \(\ell\) of
height at most \(B\). Then
\begin{equation}
 N_{\Phi,\ell}(B)\sim C_\Phi\frac{B^2}{\log B},\qquad
 C_\Phi=\frac{11}{288\sqrt3\,\pi^2}
 \prod_{p\equiv2\ (3)}\frac{1+p^{-3}}{1-p^{-3}}>0.
 \label{eq:rational-count-intro}
\end{equation}
The Euler product converges absolutely.
\end{theorem}

The proof combines character-twisted Selberg--Delange estimates with
coprimality by convolution. The congruence at \(3\) couples the numerator
and denominator and contributes an order-three character projector.
Summing the \(2\)-adic and \(3\)-adic valuation strata with their exact
projective contents then gives the displayed constant. This counts
distinct targets, since \(a\mapsto y_a\) is injective; it does not count
isomorphism classes of fiber algebras.

For comparison, retain the sufficient integer family
\begin{equation}
 \mathcal A=\{a\in\Z:a>1,\ a\equiv1\pmod9,\ a\notin\Q^{\times3},
 \ p\mid a\Rightarrow p\equiv1\pmod3\}.
 \label{eq:parameter-family}
\end{equation}
Its congruence-and-support condition is
\begin{equation}
 a\equiv1\pmod9,\qquad a\notin\Q^{\times3},\qquad
 p\mid a\Longrightarrow p\equiv1\pmod3.
 \label{eq:parameter-conditions}
\end{equation}
For \(a\in\mathcal A\) the projective coordinates
\([9:-9:32a:24a+3]\) are primitive and \(H(y_a)=32a\).

\begin{theorem}[Sufficient integer family]\label{thm:count-intro}
Writing \(M_\Phi(B)=\#\{a\in\mathcal A:H(y_a)\le B\}\), we have
\begin{equation}
 M_\Phi(B)\sim\frac{G_3(1)}{96\sqrt\pi}
 \frac{B}{\sqrt{\log B}},
 \label{eq:count-intro}
\end{equation}
where
\begin{equation}
 G_3(1)=\left[\frac23L(1,\chi_{-3})
 \prod_{p\equiv2\ (3)}(1-p^{-2})\right]^{1/2}.
 \label{eq:g3-intro}
\end{equation}
\end{theorem}

Allowing rational numerator and denominator pairs accounts for the larger
power of \(B\) in the complete line count. The geometric locus remains
one-dimensional. No classification or asymptotic is asserted for the
ambient three-dimensional target of \(\Phi\).

Section~\ref{sec:fixed-map} constructs the map and accounts for the complete
fibers. Section~\ref{sec:local-global} proves the local converse, including
the exceptional parameters. Section~\ref{sec:counting} counts the
sufficient integer families, and Section~\ref{sec:rational-counting} proves
the complete rational count in the exact target height.
Section~\ref{sec:minimality} recalls why geometric degree five is minimal.
